% !TEX root = CSL-main.tex
\subsection{Size Change Termination and Finite Power Condition}

%If formulated as in Definition \ref{def:GTC}, the Global Trace Condition (GTC) is equivalent to the algorithmic 
%criterion for Size Change Termination (SCT), introduced by Neil Jones in \cite{Lee1},
%Theorem 4. We define SCT in our formalism.

\begin{definition}[Trace, path and progress]\label{def:TracePath}
Let $\alpha \leq \omega$ be an ordinal, and $\Graph = (V, \biRel, E)$ a transition graph; then:
\begin{enumerate}[i)]
\item a {\em trace} is a map $\tau:\alpha \to X$;
\item a {\em path} is a map $\ppath:\alpha \to E$ such that for all $i + 1 < \alpha$, $\target{\ppath(i)} = \source{\ppath(i+1)}$;
\item a trace $\tau$ is {\em compatible} with a path $\ppath$ (with the same domain $\alpha$), written $\tau \models \ppath$, if 
	there is $i_0 < \alpha$ such that for all $i_0 \leq i < i+1 < \alpha$,
	if $\lbl{\ppath(i)} = (S,P)$ then $\tau(i) (S \cup P) \tau(i + 1)$, and if $\tau(i) P \tau(i + 1)$ then $i$
	is a {\em progress point} of $\tau$ in $\ppath$.
\end{enumerate}
\end{definition}

A trace $\tau:\alpha \to X$, respectively a path $\ppath:\alpha \to E$ are infinite if $\alpha = \omega$, finite otherwise.
We say that $\tau$ is a trace of $\Graph$ if $\tau \models \ppath$ for some path $\ppath$ of $\Graph$. 

%If $\tau \models\ppath$ in $\Graph$
%and $\ppath$ is the path $\edgeGstar{(S,P)}{u}{v}{\Graph}$ then we say that $\tau$ progresses in $(S,P)$ if 

\begin{definition}[Size-change Termination - SCT]
\label{def:SCT}
The transition graph $\Graph$ satisfies the {\em Size-change Termination} condition 
if for all infinite path $\ppath$ of $\Graph$ there is an infinite trace $\tau$ such that
$\tau \models \ppath$, which has infinitely many progress points.
\end{definition}


\begin{lemma}
\label{lem:SCT}
A transition graph $\Graph$ satisfies the SCT condition if and only if for all self-composable 
$\edgeGstar{(S,P)}{u}{u}{\Graph}$ such that 
$(S,P) \comp (S,P) = (S,P)$ there is some $x \in X$ such that $x P x$.
\end{lemma}

In view of Lemma \ref{lem:SCT}, SCT can be formulated in many equivalent ways, for instance, as Finite Power Condition (FPC):

\begin{definition}[Finite Power Condition - FPC]
\label{def:FPC}
$\Graph$ satisfies the {\em Finite Power Condition} if for all self-composable 
$\edgestar{(S,P)}{u}{u}$ there are some $n > 0 $ and
$x \in X$ such that $x$ progresses to $x$ in $(S,P)^n$.
\end{definition}


To prove the equivalence between SCT and FPC we use a folklore
algebraic property, that all elements of a finite monoid have some positive idempotent power. For $n > 0$, define $(S,P)^n$ as the $n$-times
self-composition of $(S,P)$.

\begin{lemma}[Idempotent Power]
\label{lem:idempotent-power}
For all $(S,P) \in \biRel^*$ there is $n > 0$ such that $(S,P)^n = (S,P)^n \comp (S,P)^n$.
%Let $M$ be a finite monoid and $x \in M$, then for some $n>0$ we have $x^nx^n=x^n$, where $x^n$ is the $n$-times product of $x$ with itself.
\end{lemma}

\begin{proof}
Since $\biRel^*$ is finite, there are only finitely many powers of $(S,P)$, therefore 
for some $a,b>0$ we have $(S,P)^a = (S,P)^{a+b}$. Then for all $c \ge a$ we have $(S,P)^c = (S,P)^{c+b}$,
hence $(S,P)^c = (S,P)^{c+db}$ for all $d>0$. Take the first $d>0$ such that $db \ge a$: then 
$(S,P)^{db} = (S,P)^{db+db} = (S,P)^{db} \comp (S,P)^{db}$. We take $n=db$.
\end{proof}

%\begin{remark}\label{rem:idempotent-power}
%Now $\biRel^*$, which is finite, but neither $\biRel^*$ nor the set of bi-partite relations that are composable in $\Graph$ is a monoid because of the lack of identities, hence of the unit w.r.t. the composition $\comp$. 
%Nonetheless, the proof of the above proposition does not rely on the existence of the unit in $M$, only on the associativity of its product, hence it applies
%to $(\biRel^*, \comp)$.
%\end{remark}

\begin{proposition}[FPC and SCT are equivalent]
\label{proposition-SCT}
SCT $\Leftrightarrow$ FPC
\end{proposition}

\begin{proof}
%\begin{description}
%\item
Assume SCT, $\edgeGstar{(S,P)}{u}{u}{\Graph}$.
By Lemma \ref{lem:idempotent-power} and finiteness of the bi-partite graph $\Graph$
there is some $n > 0$ such that $(S,P)^n \comp (S,P)^n=(S,P)^n$. By SCT,
there is some $x \in X$ such that $x$ progresses to $x$ in $(S,P)^n$.
%\item

Assume FPC, $\edgeGstar{(S,P)}{u}{u}{\Graph}$ and 
$(S,P) \comp (S,P)=(S,P)$ in order to prove that 
there is some $x \in X$ such that $x$ progresses to $x$ in $(S,P)$.
By FPC, there are some $n > 0 $ and $x \in X$ such that $x$ progresses to $x$ in 
$(S,P)^n$. By $(S,P)\comp (S,P)=(S,P)$ we deduce that 
$(S,P)^n=(S,P)$, hence $x$ progresses to $x$ in $(S,P)$.
%\end{description}
\end{proof}


